\part{Chimie (7 points)~: L'acide éthanoïque et ses utilisations}
L'acide éthanoïque de formule $\ce{CH3COOH}$ représente le principal constituant
du vinaigre commercial après l'eau. Il est utilisé comme réactif dans de
nombreuses synthèses organiques comme celle qui conduit à l'éthanoate
d'éthyle.

Le degré d'acidité d'un vinaigre est donné en degré ($^{\circ}$).

Cet exercice se compose de 3 parties indépendantes et vise~:
\begin{itemize}
    \item l'étude d'une solution aqueuse d'acide éthanoïque~;
    \item la détermination du degré d'acidité (titre) d'un vinaigre commercial~;
    \item l'étude de la synthèse de l'éthanoate d'éthyle à partir de l'acide
          éthanoïque.
\end{itemize}

\begin{donnees}
    \begin{itemize}
        \item Le degré d'acidité d'un vinaigre est égal à la masse, en grammes
              d'acide pur contenue dans $\qty{100}{\mL}$ de ce vinaigre.
        \item $\mathrm{pK_{A}}
              (\ce{CH3COOH_{(aq)}}\ttslash{}\ce{CH3COO^-_{(aq)}})=4,8$ à
              $\qty{25}{\degreeCelsius}$~; \;
              $M(\ce{CH3COOH})=\qty{60}{\g\per\mol}$.
    \end{itemize}
\end{donnees}

\section*{Partie 1~: Étude d'une solution aqueuse d'acide éthanoïque}
La mesure du $\mathrm{pH}$ d'une solution aqueuse d'acide éthano\"ique, à
$\qty{25}{\degreeCelsius}$, a donné $\mathrm{pH}=3,0$.

\begin{enumerate}
    \item Écrire l'équation chimique modélisant la transformation entre l'acide
          éthanoïque et l'eau.
    \item Déterminer l'espèce du couple
          $(\ce{CH3COOH_{(aq)}}\ttslash{}\ce{CH3COO^-_{(aq)}})$ qui prédomine
          dans la solution. Justifier.
    \item Déterminer la valeur du quotient de réaction $Q_{r,\eq}$ à l'état
          d'équilibre du système chimique.
    \item La valeur de $Q_{r,\eq}$ est-elle modifiée si on dilue la solution
          d'acide éthanoïque~? Justifier.
\end{enumerate}



\section*{Partie 2~: Détermination du degré d'acidité d'un vinaigre commercial}
L'étiquette d'une bouteille de vinaigre commercial indique $6^{\circ}$. La
concentration molaire en acide éthanoïque dans ce vinaigre est $C_{0}$.

On se propose de doser par pH-métrie ce vinaigre afin de déterminer son degré
d'acidité. Pour cela, on prépare une solution aqueuse $\mathrm{(S_{1})}$ par
dilution 10 fois du vinaigre commercial et on prélève un volume
${V_{A}=\qty{25}{\mL}}$ de la solution diluée $\mathrm{(S_{1})}$ de concentration
molaire $C_{A}$ $\left(C_{A}=\frac{C_{0}}{10}\right)$ que l'on dose avec une
solution aqueuse $\mathrm{(S_{2})}$ d'hydroxyde de sodium
$\ce{Na^{+}_{(aq)} + HO^{-}_{(aq)}}$ de concentration molaire
$C_{B}=\qty{2.5e-1}{\mol\per\L}$.

À l'équivalence, le volume de la solution $\mathrm{(S_{2})}$ ajouté est
$V_{B,E}=\qty{10}{\mL}$.

\begin{enumerate}
    \item Écrire l'équation de la réaction qui a eu lieu lors du dosage,
          supposée totale.
    \item Calculer la valeur de $C_{A}$. En déduire la valeur de $C_{0}$.
    \item Vérifier la valeur du degré d'acidité du vinaigre indiquée sur
          l'étiquette de la bouteille.
\end{enumerate}


\section*{Partie 3~: Synthèse de l'éthanoate d'éthyle à partir de l'acide
éthanoïque}
Dans un ballon, on introduit un mélange équimolaire de $n_{1}=0,3 \mathrm{~mol}$ d'acide éthanoïque et $n_{2}=0,3 \mathrm{~mol}$ d'éthanol et quelques gouttes d'acide sulfurique concentré. À l'état d'équilibre du système chimique, la quantité de matière d'ester formé est~: $n_{f}(e s t e r)=0,2 \mathrm{~mol}$. La synthèse de l'éthanoate d'éthyle est modélisée par la réaction d'équation~:

$$
CH3 \mathrm{COOH}_{(\ell)}+C2 H5 \mathrm{OH}_{(\ell)} \rightleftharpoons CH3 COOC2 \mathrm{H}_{5(\ell)}+H2 \mathrm{O}_{(\ell)}
$$

\begin{enumerate}
  \item Identifier les groupes caractéristiques des molécules organiques figurant dans l'équation de la réaction de synthèse.
  \item Donner les caractéristiques de cette réaction.
  \item Déterminer la valeur du rendement de cette synthèse.
  \item Déterminer la valeur de la constante d'équilibre K associée à l'équation chimique de la réaction d'estérification.
  \item pour synthétiser l'éthanoate d'éthyle par une transformation rapide et totale, il est possible de remplacer l'acide éthanoïque par l'un de ses dérivés.\\
Donner la formule semi-développée de ce dérivé.
\end{enumerate}

